\section{Methodology}

In this work, TALYS simulations were compared with experimental outgoing-particle energy spectra from EXFOR. The main objective was to evaluate how well TALYS reproduces the measured light-ion spectra and to investigate the sensitivity of the results to different pre-equilibrium model settings. The reactions studied were:

\begin{itemize}
\item proton production, ((n,px)), on selected targets, including \ce{^{28}Si}, \ce{^{56}Fe}, \ce{^{52}Cr}, \ce{^{51}V}, and \ce{^{59}Co};
\item deuteron production, ((n,dx)), on \ce{^{28}Si};
\item alpha production, ((n,\alpha x)), on \ce{^{28}Si}.
\end{itemize}

For each reaction case, TALYS calculations were performed for the incident neutron energies available in the selected EXFOR datasets, or for a representative subset of those energies when the dataset contained a large number of points. The outgoing-particle energy is denoted by (E_{\mathrm{out}}), and the calculated spectra were compared with the corresponding EXFOR energy-differential cross sections.

To study the sensitivity of the calculated spectra to the treatment of pre-equilibrium emission, four TALYS calculations were performed for each reaction case, corresponding to \texttt{preeqmode = 1--4}. According to the TALYS manual, these options select different models for pre-equilibrium reactions, ranging from exciton-model calculations with different treatments of transition rates to a multi-step direct/compound model \cite{koning2025talys}. In this work, the four calculations are labelled as \texttt{preeq1}, \texttt{preeq2}, \texttt{preeq3}, and \texttt{preeq4}. All other input conditions were kept fixed, so that differences between the spectra could be attributed mainly to the selected pre-equilibrium model.

\begin{table}[H]
\centering
\caption{TALYS pre-equilibrium models compared in this work. Definitions follow the TALYS manual \cite{koning2025talys}.}
\label{tab:preeq-configurations}
\begin{tabular}{c c p{0.62\textwidth}}
\hline
\textbf{Label used in this work} & \textbf{TALYS keyword} & \textbf{Model definition} \
\hline
\texttt{preeq1} & \texttt{preeqmode 1} & Exciton model: analytical transition rates with energy-dependent matrix element. \

\texttt{preeq2} & \texttt{preeqmode 2} & Exciton model: numerical transition rates with energy-dependent matrix element. \

\texttt{preeq3} & \texttt{preeqmode 3} & Exciton model: numerical transition rates with optical model for collision probability. \

\texttt{preeq4} & \texttt{preeqmode 4} & Multi-step direct/compound model. \
\hline
\end{tabular}
\end{table}


The comparison was based on both qualitative and quantitative criteria. First, the TALYS total spectra were plotted together with the EXFOR data in linear scale, in order to compare the overall shape, the peak position, and the peak magnitude. The same comparison was also shown in logarithmic scale, which makes it easier to inspect the low-cross-section region and the high-(E_{\mathrm{out}}) tail.

Second, the TALYS spectra were compared with EXFOR through a point-by-point ratio. Since the TALYS output and the EXFOR data are not always given at exactly the same outgoing energies, the TALYS total spectrum was interpolated at the EXFOR (E_{\mathrm{out}}) values. The ratio was then calculated as

\begin{equation}
R(E_{\mathrm{out}}) =
\frac{\left(\mathrm{d}\sigma/\mathrm{d}E\right)*{\mathrm{TALYS}}(E*{\mathrm{out}})}
{\left(\mathrm{d}\sigma/\mathrm{d}E\right)*{\mathrm{EXFOR}}(E*{\mathrm{out}})} .
\label{eq:ratio}
\end{equation}

A value (R(E_{\mathrm{out}}) \simeq 1) indicates good agreement between TALYS and EXFOR at that outgoing energy. Values above 1 indicate that TALYS overpredicts the experimental data, while values below 1 indicate an underprediction. The ratio plot is especially useful when the total spectra appear visually similar, because it makes systematic deviations easier to identify.

Finally, for the selected pre-equilibrium mode in each case, the TALYS spectrum was decomposed into its main reaction-mechanism components. This made it possible to identify which parts of the spectrum are mainly associated with compound emission, pre-equilibrium emission, direct-like processes, or multiple pre-equilibrium emission. In general, the low-(E_{\mathrm{out}}) peak is expected to be more strongly associated with compound-nucleus evaporation, while the shoulder and high-(E_{\mathrm{out}}) tail are more sensitive to pre-equilibrium and direct-like contributions.

Each diagnostic figure therefore contains four panels: the total TALYS--EXFOR comparison in linear scale, the same comparison in logarithmic scale, the ratio (R(E_{\mathrm{out}})) for all four \texttt{preeqmode} options, and the mechanism decomposition for the selected mode. The best-performing mode was chosen by considering the total spectrum, the ratio behaviour, and the physical interpretation of the component decomposition. Particular attention was given to the shoulder and high-energy tail, because these regions are often more sensitive to the pre-equilibrium treatment than the main evaporation peak.

In the logarithmic plots, very small TALYS values near the numerical or physical end of the spectrum were interpreted with caution. These regions can visually exaggerate differences between modes, especially when EXFOR has few or no data points there. For this reason, the final assessment focused mainly on the outgoing-energy interval where experimental data are available and physically meaningful comparisons can be made.
